\section*{§9. Распределение Больцмана}
\addcontentsline{toc}{section}{§9. Распределение Больцмана}

\begin{basics}
\sub{Распределение Больцмана}
\[n(\vec r)=n_0\,e^{-U(\vec r)/kT},\qquad
w(\vec r)=\frac{e^{-U(\vec r)/kT}}{\int_Ve^{-U/kT}\dd^3r},\qquad n(\vec r)=Nw(\vec r).\]
\begin{itemize}
\item $U(\vec r)$ — \emph{потенциальная энергия одной молекулы} во внешнем поле (Дж), а не потенциал поля.
Задана с точностью до константы (она сокращается). Если дана сила: $\vec F=-\nabla U$.
\item $n(\vec r)$ — концентрация (молекул на м$^3$); $n_0$ — концентрация там, где $U=0$.
\item $w(\vec r)$ — плотность вероятности найти \emph{данную} молекулу в точке $\vec r$: $w\,\dd^3r$ — вероятность
оказаться в объёме $\dd^3r$. Нормировка $\int_Vw\,\dd^3r=1$, т.е.\ $\int_Vn\,\dd^3r=N$ — отсюда находят $n_0$.
\item Применимость: идеальный газ (взаимодействие молекул мало), равновесие, $T$ одинакова везде.
\item Давление: $p(\vec r)=n(\vec r)kT$ — меняется так же, как $n$.
\end{itemize}

\sub{Не путать плотности}
\begin{tabular}{@{}lll@{}}
$w(\vec r)$ & на единицу объёма & $\propto e^{-U/kT}$\\
$w(z)$ & по высоте & $\propto S(z)e^{-U/kT}$ ($S$ — сечение)\\
$w(r)$ (шар) & по расстоянию от центра & $\propto4\pi r^2e^{-U/kT}$\\
$w(r)$ (цилиндр) & по расстоянию от оси & $\propto2\pi r\,e^{-U/kT}$
\end{tabular}
\\«Наиболее вероятное расстояние» — максимум $w(r)$ \emph{с множителем} $r^2$ или $r$.

\sub{Поле тяжести}
$U=mgz$ ($z$ — высота вверх):
\[n(z)=n(0)e^{-mgz/kT},\qquad p(z)=p(0)e^{-\mu gz/RT}\ \text{(барометрическая формула)}.\]
Полубесконечный цилиндр: $w(z)=\frac{mg}{kT}e^{-mgz/kT}$, $\avg z=\frac{kT}{mg}$, $\avg{z^2}=2\big(\frac{kT}{mg}\big)^2$, $\avg{mgz}=kT$.
\\Столб высоты $H$, $x=mgH/kT$: $\avg{mgz}=kT-\dfrac{mgH}{e^{x}-1}$.
\\Центр масс: $z_c=\dfrac{\sum m_i\avg{z_i}}{\sum m_i}$, для одного газа $z_c=\avg z$.

\sub{Вращающийся сосуд (центрифуга)}
В системе, вращающейся с $\omega$, на молекулу действует центробежная сила $m\omega^2r$, ей отвечает
$U=-\frac{m\omega^2r^2}{2}$ ($r$ — расстояние от оси):
$n(r)=n(0)\,e^{m\omega^2r^2/2kT}$ — газ прижимается к стенке, тяжёлые молекулы сильнее.

\sub{Несколько полей}
$U=U_1+U_2$ $\Rightarrow$ $e^{-U/kT}=e^{-U_1/kT}e^{-U_2/kT}$: координаты независимы, средние энергии складываются.
Полная энергия молекулы $\avg E=\frac32kT+\avg U$.
\end{basics}

%-------------------------------------------------------------
\begin{problem}{9.1}{exdz}{Найти среднее значение и относительную флуктуацию потенциальной энергии молекул идеального газа,
находящегося в прямом полубесконечном вертикальном цилиндре в однородном поле тяжести. Температура $T$, масса молекулы $m$.}
\textbf{Найти:} $\avg U$ и $\delta(U)=\sigma(U)/\avg U$, где $U=mgz$. Ось $z$ — от дна вверх, $z\in[0,\infty)$; $a=kT/mg$.

\textbf{1. Плотность вероятности по высоте.} По Больцману $w(z)=Ce^{-mgz/kT}=Ce^{-z/a}$ (сечение постоянное и сокращается).
Нормировка: $\int_0^\infty Ce^{-z/a}\dd z=Ca=1\Rightarrow w(z)=\frac1ae^{-z/a}$.

\textbf{2. Средние $z$ и $z^2$.} По формуле $\int_0^\infty z^ke^{-z/a}\dd z=k!\,a^{k+1}$:
\[\avg z=\frac1a\cdot a^2=a,\qquad \avg{z^2}=\frac1a\cdot2a^3=2a^2.\]

\textbf{3. Среднее и дисперсия $U=mgz$.} Множитель $mg$ выносится:
\[\avg U=mg\,\avg z=mga=kT,\qquad
\sigma^2(U)=\avg{U^2}-\avg U^2=(mg)^2\big(\avg{z^2}-\avg z^2\big)=(mg)^2a^2=(kT)^2.\]

\textbf{4. Относительная флуктуация:} $\delta(U)=\sigma(U)/\avg U=kT/kT=1$.
\answer{$\avg U=kT$, $\delta(U)=1$.}
\note{разброс энергии одной молекулы равен её среднему значению (100\%).}
\end{problem}

%-------------------------------------------------------------
\begin{problem}{9.2}{exdz}{Идеальный газ из $N$ молекул массой $m$ находится в вертикальном цилиндре высотой $H$ с площадью основания $\pi R^2$.
Температура $T$. Найти среднюю потенциальную энергию молекулы $\avg U$, высоту центра масс газа $z_c$ и во сколько раз
концентрация на высоте $z_c$ меньше, чем у дна: $\alpha=n(z_c)/n(0)$.}
Ось $z$ — от дна вверх. Обозначим $a=kT/mg$.

\textbf{1. Плотность вероятности по высоте.} По Больцману $w(z)=Ce^{-mgz/kT}=Ce^{-z/a}$ (сечение постоянное и сокращается).
Молекулы есть только при $0\le z\le H$, поэтому нормируем на этом отрезке:
\[\int_0^HCe^{-z/a}\dd z=Ca\big(1-e^{-H/a}\big)=1\ \Rightarrow\
w(z)=\frac{mg}{kT}\,\frac{e^{-mgz/kT}}{1-e^{-mgH/kT}}.\]

\textbf{2. Средняя высота.} $\avg z=\int_0^Hz\,w(z)\dd z$. Интеграл берём по частям:
\[\int_0^Hze^{-z/a}\dd z=\Big[-aze^{-z/a}\Big]_0^H+a\int_0^He^{-z/a}\dd z=-aHe^{-H/a}+a^2\big(1-e^{-H/a}\big).\]
Делим на нормировку $a(1-e^{-H/a})$:\quad $\avg z=a-\dfrac{H}{e^{H/a}-1}$.

\textbf{3. Ответы.} $\avg U=mg\avg z$; центр масс одинаковых молекул — их средняя высота, $z_c=\avg z$;
концентрация $n\propto e^{-mgz/kT}$, поэтому $\alpha=e^{-mgz_c/kT}$.
\answer{$\avg U=kT-\dfrac{mgH}{e^{mgH/kT}-1}$,\; $z_c=\dfrac{kT}{mg}-\dfrac{H}{e^{mgH/kT}-1}$,\; $\alpha=e^{-mgz_c/kT}$.}
\note{при $H\to\infty$ получается 9.1: $\avg U=kT$, $z_c=kT/mg$.}
\end{problem}

%-------------------------------------------------------------
\begin{problem}{9.3}{exdz}{Пылинки массой $m_0=10^{-21}$\,кг взвешены в воздухе при $T=300$\,К. Оценить толщину $\Delta z$ слоя воздуха у
поверхности Земли, в пределах которого концентрация пылинок отличается не более чем на 1\%.}
\textbf{Найти:} $\Delta z$, при котором $\Delta n/n=0{,}01$.

\textbf{1.} Пылинки ведут себя как молекулы газа в поле тяжести — распределение Больцмана:
$n(z)=n(0)\,e^{-m_0gz/kT}$.

\textbf{2.} Слой тонкий, поэтому изменение концентрации $\approx$ производная $\times$ толщина:
\[\Delta n\approx\Big|\frac{\dd n}{\dd z}\Big|\Delta z=\frac{m_0g}{kT}\,n\,\Delta z\quad\Rightarrow\quad
\frac{\Delta n}{n}=\frac{m_0g}{kT}\Delta z.\]

\textbf{3.} Выражаем $\Delta z$ и подставляем числа: $kT=1{,}38\cdot10^{-23}\cdot300=4{,}14\cdot10^{-21}$\,Дж,
$m_0g=9{,}8\cdot10^{-21}$\,Н:
\[\Delta z=\frac{kT}{m_0g}\cdot\frac{\Delta n}{n}=0{,}42\text{ м}\cdot0{,}01.\]
\answer{$\Delta z\approx4{,}2$ мм.}
\end{problem}

%-------------------------------------------------------------
\begin{problem}{9.4}{exdz}{В центрифуге радиуса $R$ и высоты $H$, вращающейся с угловой скоростью $\omega$, разделяется смесь двух газов
с массами молекул $m_1$ и $m_2$. Найти коэффициент разделения $q=\dfrac{n_1(R)/n_1(0)}{n_2(R)/n_2(0)}$, где $n_{1,2}(r)$ — концентрации
газов на расстоянии $r$ от оси. Силу тяжести не учитывать, температура $T$.}
\textbf{1. Потенциальная энергия.} Во вращающейся системе на молекулу действует центробежная сила $F=m_i\omega^2r$ (от оси).
Ей соответствует $U_i=-\int F\,\dd r=-\dfrac{m_i\omega^2r^2}{2}$.

\textbf{2. Больцман для каждого газа отдельно} (газы идеальные, не влияют друг на друга):
\[n_i(r)=n_i(0)\,e^{-U_i/kT}=n_i(0)\,e^{m_i\omega^2r^2/2kT}\ \Rightarrow\ \frac{n_i(R)}{n_i(0)}=e^{m_i\omega^2R^2/2kT}.\]

\textbf{3. Делим} выражение для газа 1 на выражение для газа 2.
\answer{$q=\exp\Big[\dfrac{(m_1-m_2)\omega^2R^2}{2kT}\Big]$.}
\note{UF$_6$ с $^{235}$U и $^{238}$U, $R=1$\,м, $T=300$\,К, $\omega=1300$\,с$^{-1}$: $|m_1-m_2|=3\cdot1{,}66\cdot10^{-27}$\,кг, показатель $\approx1{,}0$;
в книге $q\approx2{,}7$ — это если тяжёлый изотоп поставить в числитель.}
\end{problem}

%-------------------------------------------------------------
\begin{problem}{9.5}{ex}{Смесь двух идеальных газов: $N_1$ частиц массой $m_1$ и $N_2$ частиц массой $m_2$ — находится в прямом
цилиндрическом сосуде (высоту считаем очень большой). Найти положение центра масс смеси, если температура $T$, $\vec g$ направлено вдоль оси.}
\textbf{1. Центр масс} — по определению среднее высот, взвешенное по массам:
\[z_c=\frac{N_1m_1\avg{z_1}+N_2m_2\avg{z_2}}{N_1m_1+N_2m_2}.\]

\textbf{2. Средняя высота молекулы каждого газа} — из 9.1 (газы независимы): $\avg{z_i}=\dfrac{kT}{m_ig}$, поэтому $m_i\avg{z_i}=\dfrac{kT}{g}$ для обоих.
\answer{$z_c=\dfrac{(N_1+N_2)\,kT}{g\,(N_1m_1+N_2m_2)}$.}
\note{при конечной высоте $H$ брать $\avg{z_i}$ из 9.2.}
\end{problem}

%-------------------------------------------------------------
\begin{problem}{9.8}{dz}{Идеальный газ из $N$ частиц находится в сосуде в виде перевёрнутого прямого конуса (вершина внизу, на поверхности Земли),
ось параллельна $\vec g$, угол при вершине $2\beta$, высота бесконечно большая. Найти концентрацию $n_0$ молекул у вершины. Температура $T$, масса молекулы $m$.}
\textbf{1. Больцман:} $n(z)=n_0e^{-mgz/kT}$, $z$ — высота над вершиной.

\textbf{2. Объём тонкого слоя.} На высоте $z$ сечение — круг радиуса $z\tan\beta$, площадь $S(z)=\pi z^2\tan^2\beta$, объём слоя $S(z)\dd z$.

\textbf{3. Условие нормировки:} всех молекул $N$:
\[N=\int_0^\infty n(z)S(z)\dd z=n_0\pi\tan^2\beta\int_0^\infty z^2e^{-mgz/kT}\dd z=n_0\pi\tan^2\beta\cdot2\Big(\frac{kT}{mg}\Big)^3\]
(использовали $\int_0^\infty z^2e^{-z/a}\dd z=2a^3$, $a=kT/mg$). Выражаем $n_0$.
\answer{$n_0=\dfrac{N(mg)^3}{2\pi\tan^2\beta\,(kT)^3}$.}
\note{так же 9.7 (цилиндр с площадью дна $S$): $N=n_0S\frac{kT}{mg}\Rightarrow n_0=\frac{Nmg}{SkT}$.}
\end{problem}

%-------------------------------------------------------------
\begin{problem}{9.12}{dz}{Отрицательные ионы массой $m$ и зарядом $q$ находятся при температуре $T$ между пластинами плоского конденсатора.
Поле $\vec E$ направлено туда же, куда $\vec g$ (вниз). Во сколько раз концентрация ионов у верхней пластины отличается от
концентрации у нижней, если расстояние между пластинами $d$?}
\textbf{Найти:} $n_{\text{верх}}/n_{\text{низ}}$. Ось $z$ — вверх, $z=0$ у нижней пластины.

\textbf{1. Силы.} Тяжесть $mg$ — вниз. Ион отрицательный, поэтому электрическая сила $|q|E$ направлена \emph{против} $\vec E$, т.е.\ вверх.

\textbf{2. Потенциальная энергия} (сила вниз $\Rightarrow U$ растёт с высотой):
$U(z)=mgz-|q|Ez=(mg-|q|E)z$.

\textbf{3. Больцман:} $\dfrac{n(d)}{n(0)}=e^{-[U(d)-U(0)]/kT}=e^{-(mg-|q|E)d/kT}$.
\answer{$\dfrac{n_{\text{верх}}}{n_{\text{низ}}}=\exp\Big[\dfrac{d\,(|q|E-mg)}{kT}\Big]$.}
\end{problem}

%-------------------------------------------------------------
\begin{problem}{9.13}{dz}{Идеальный газ находится в поле центральной силы с потенциальной энергией $U(r)=\alpha r^2$ ($\alpha>0$, $r$ — расстояние до центра).
Температура $T$, число частиц $N$. Найти $n(r)$ и наиболее вероятное расстояние $r_0$ молекул от центра. Найти плотность вероятности $w(U)$
и наиболее вероятное значение потенциальной энергии $\tilde U$.}
\textbf{1. Концентрация.} По Больцману $n(r)=n_0e^{-\alpha r^2/kT}$. Константу $n_0$ — из нормировки по всему пространству (слой $4\pi r^2\dd r$):
\[N=n_0\int_0^\infty4\pi r^2e^{-\alpha r^2/kT}\dd r=n_0\cdot4\pi\cdot\frac{\sqrt\pi}{4}\Big(\frac{kT}{\alpha}\Big)^{3/2}=n_0\Big(\frac{\pi kT}{\alpha}\Big)^{3/2}\]
(гауссов интеграл $I_2=\frac{\sqrt\pi}{4a^{3/2}}$ с $a=\alpha/kT$).

\textbf{2. Наиболее вероятное расстояние.} Вероятность оказаться в слое $(r,r+\dd r)$ пропорциональна \emph{объёму слоя}:
$w(r)\propto r^2e^{-\alpha r^2/kT}$. Максимум: $\big(2r-\frac{2\alpha r^3}{kT}\big)e^{\dots}=0\Rightarrow r_0=\sqrt{kT/\alpha}$.

\textbf{3. Распределение по $U$.} Замена $U=\alpha r^2$: $r=\sqrt{U/\alpha}$, $\dd r=\frac{\dd U}{2\sqrt{\alpha U}}$:
\[w(U)\propto r^2e^{-U/kT}\frac{\dd r}{\dd U}\propto\frac U\alpha\cdot\frac{1}{\sqrt{\alpha U}}e^{-U/kT}\propto\sqrt U\,e^{-U/kT}.\]
Это та же форма, что $w(E)$ в 8.2, поэтому нормировочная константа та же. Максимум: $\frac{\dd}{\dd U}\big(\sqrt Ue^{-U/kT}\big)=0\Rightarrow\tilde U=\frac{kT}2$.
\answer{$n(r)=N\Big(\dfrac{\alpha}{\pi kT}\Big)^{3/2}e^{-\alpha r^2/kT}$,\; $r_0=\sqrt{kT/\alpha}$,\;
$w(U)=\dfrac{2}{\sqrt\pi}(kT)^{-3/2}\sqrt U\,e^{-U/kT}$,\; $\tilde U=\dfrac{kT}2$.}
\note{$\tilde U\ne\alpha r_0^2=kT$: наиболее вероятные значения не пересчитываются через замену.}
\end{problem}

%-------------------------------------------------------------
\begin{problem}{9.14}{dz}{Газ находится в полубесконечном прямом цилиндре в однородном поле тяжести. Найти долю молекул, потенциальная
энергия которых больше их средней кинетической энергии.}
\textbf{Найти:} долю молекул с $mgz>\avg{E_{\text{кин}}}=\frac32kT$, т.е.\ выше высоты $z_0=\frac{3kT}{2mg}$.

\textbf{1.} Плотность по высоте (9.1): $w(z)=\frac{mg}{kT}e^{-mgz/kT}$.

\textbf{2.} Доля выше $z_0$:
$\displaystyle P=\int_{z_0}^\infty w(z)\dd z=e^{-mgz_0/kT}=e^{-3/2}$.
\answer{$e^{-1{,}5}\approx0{,}22$.}
\end{problem}

%-------------------------------------------------------------
\begin{problem}{9.15}{dz}{В центрифуге радиуса $R$ и высоты $H$, вращающейся с угловой скоростью $\omega$, при температуре $T$ находится $N$ молекул
идеального газа массой $m$. Найти среднее значение полной энергии молекулы. Ускорение $\vec g$ направлено вдоль оси.}
\textbf{1. Из чего состоит энергия.} $E=E_{\text{кин}}+U_{\text{ц}}+U_g$, где $U_{\text{ц}}=-\frac{m\omega^2r^2}2$ (центробежная), $U_g=mgz$.
Среднее суммы = сумма средних. Кроме того, $e^{-(U_{\text{ц}}+U_g)/kT}=e^{-U_{\text{ц}}/kT}e^{-U_g/kT}$ — по $r$ и по $z$ распределения независимы,
каждое слагаемое считаем отдельно. $\avg{E_{\text{кин}}}=\frac32kT$.

\textbf{2. Тяжесть} — ответ 9.2: $\avg{U_g}=kT-\dfrac{mgH}{e^{mgH/kT}-1}$.

\textbf{3. Центробежная часть.} По радиусу: $w(r)\propto2\pi r\,e^{m\omega^2r^2/2kT}$ (площадь кольца $2\pi r\,\dd r$).
Удобная переменная $s=\frac{m\omega^2r^2}{2kT}\in[0,b]$, $b=\frac{m\omega^2R^2}{2kT}$; так как $\dd s\propto r\,\dd r$, то $w(s)=\frac{e^s}{e^b-1}$. Тогда
\[\avg s=\frac{\int_0^bse^s\dd s}{e^b-1}=\frac{(b-1)e^b+1}{e^b-1}=\frac{b\,e^b}{e^b-1}-1,\qquad
\avg{U_{\text{ц}}}=-kT\avg s=kT-\frac{m\omega^2R^2}{2(1-e^{-b})}\]
(это ответ 9.9).

\textbf{4. Складываем:} $\frac32kT+kT+kT=\frac72kT$ плюс две поправки.
\answer{$\avg E=\dfrac{7kT}2-\dfrac{m\omega^2R^2}{2\big[1-e^{-m\omega^2R^2/2kT}\big]}-\dfrac{mgH}{e^{mgH/kT}-1}$.}
\end{problem}

%-------------------------------------------------------------
\begin{problem}{9.16}{dz}{Два прямых цилиндрических сосуда высотой $H$ каждый стоят друг на друге и содержат по одному молю идеального газа при
температуре $T$. На какой высоте от поверхности Земли находится центр масс системы? Масса молекулы $m$.}
\textbf{1. Нижний сосуд} ($0\le z\le H$) — ответ 9.2: $z_1=\dfrac{kT}{mg}-\dfrac{H}{e^{x}-1}$, $x=\dfrac{mgH}{kT}$.

\textbf{2. Верхний сосуд} — то же распределение, но отсчитанное от его дна на высоте $H$: $z_2=z_1+H$.

\textbf{3. Массы газов одинаковы} (по молю), поэтому центр масс — посередине: $z_c=\frac{z_1+z_2}{2}=z_1+\frac H2$.
\answer{$z_c=\dfrac{kT}{mg}+\dfrac H2-\dfrac{H}{e^{mgH/kT}-1}$.}
\end{problem}
